\section{Global native entry from actual physical contacts}
\label{sec:global-native-entry-rederived}

This section replaces a source-chart assumption by an actual implication.
It uses the full-parent dictionary and SE1--SE10, whose proofs establish
contact coercivity, strict source-field curvature, contact sufficiency,
and the event-deletion test.  Those are single-field assertions; no
strict residue or mixed tensor inverse is inferred here.

\subsection{Statement and the real maximal-contact graph}

Use the raw full parent \(F\), with \(d>0\).
Every source below is physical and has an actual finite noncommuting
zero.  The two allowed coefficient domains are
\[
\begin{array}{ll}
 {\rm C1}:&\lambda=n=0,\quad m,C,A,\alpha>0,\quad
 D,\beta\ge0,\quad D+\beta>0,\\
 {\rm C2}:&\lambda=\beta=0,\quad m,C,A,D,\alpha>0,\quad n\ge0.
\end{array}\tag{EN1}
\]
Make the actual angle gauges positive and put \(Z_c=1\), as permitted
by SE2.  Then \(W_c>0\) by SE9.  We prove the strict signs
\[
\begin{gathered}
 X_c<0,\quad Y_c>0,\quad u_cW_c-Y_c>0,\quad
 1-v_cY_c>0,\\
 R_{m,c}:=v_cW_c-X_c+u_cv_cY_c-u_c>0,\\
 R_{D,c}:=u_cW_c+u_cv_cX_c-Y_c-v_c<0,\\
 X_c+u_c>0,\qquad W_c+v_cX_c>0.
\end{gathered}\tag{EN2}
\]
These inequalities hold over every real closed coefficient field.
In particular they are strict ordered-field statements, with no
assertion that a gap has positive real residue.

Here is the common connected-graph argument, proved over \(\mathbb R\)
before first-order transfer.  In either exact two-omission family,
initially require all six retained events positive, normalize \(C=2\),
and let their other five weights vary over the positive orthant.  Put
\[
 d_*(w)=\max\{d\ge0:(d,w)\text{ is physical}\}.
 \tag{EN3}
\]
The admissible set is a closed interval by convexity.  It has positive
length: at \(d=0\) the central Laurent coefficient is positive, its
nonnegative Cauchy--Binet expansion containing \(16mCAD\); the outer
coefficients remain nonnegative for \(d\ge0\).  Thus sufficiently small
positive \(d\) is physical by SE1--SE2 and the homogeneous endpoint
limits.  A fixed strictly CHSH-violating real test bounds \(d\) above.

For these positive six-event supports the active outer exponents are
constant and have zero in the interior of their convex hull.  The
logarithmic Laurent sum is strictly convex and coercive.  At \(d_*\)
its unique minimum is zero.  Indeed if its minimum were \(\mu>0\),
a small increase of \(d\) preserves every active outer coefficient
within a relative error \(\varepsilon\) and the central coefficient
within absolute error \(\varepsilon\).  As in SE13, this preserves a
positive uniform lower bound, contradicting maximality.
There is therefore one positive squared-angle contact, and SE2 gives
one state with \(Z=1\).

At that contact the core value is strictly negative (SE8--SE9).
The Jacobian of the six equations
\(\nabla_zF=0,F=0\), in \(z=(W,X,Y,u,v)\) and \(d\), is
\[
 \begin{pmatrix}\nabla_z^2F&\partial_d\nabla_zF\\
 0&\partial_dF\end{pmatrix},
 \qquad \nabla_z^2F\succ0,\quad\partial_dF<0 .
 \tag{EN4}
\]
The real analytic implicit-function theorem gives a local analytic
contact branch.  SE7 makes it physical and the negative core value
makes its \(d\) maximal.  Uniqueness makes these branches agree.
Thus the normalized contact is analytic on the entire connected
positive orthant.  Every physical finite zero with those event weights
lies on this graph, because its negative core value prevents increasing
\(d\) at that same test.

\subsection{C1: strict signs including the two literal boundary faces}

Write
\[
 (W,X,Y,Z)=(S,u(\theta-1),uY_0,1).
 \tag{EN5}
\]
The following two identities are the actual retained sign-entry
identities, also checked against the full raw parent in the focused
checker.  With \(B=\operatorname{adj}P\),
\[
\begin{gathered}
 uB_{44}+B_{24}
 =16uv^2D_u^2D_v\,\mathcal F,\\
\begin{aligned}
 \mathcal F={}&(d+2D)\{(d+2\alpha)(mD_v+Cu^2v^2)
                       +2AD_v(m+Cu^2)\}\\
 &+2mC\{d(u^2v^2+D_v)+2\alpha+2Dv^2D_u\}>0 .
\end{aligned}
\end{gathered}\tag{EN6}
\]
Since \(P\) has one-dimensional kernel and three positive eigenvalues,
\(B=\gamma\psi\psi^T\) with \(\gamma>0\).  Equation EN6 gives
\(\theta>0\).

Neither \(Y_0=0\) nor \(S-Y_0=0\) is possible.  These walls would delete
the positive \(A\)-event or \(C\)-event, respectively, while preserving
the kernel and both angle jets.  The resulting omission set contains
an exclusive pair, so has noncoercive support.  SE7 and SE10 rule this out.
The remaining cross-multiplied contact identity is
\[
 m\,[S+uv(\theta-1)]\,
       [v(S+u^2Y_0)-u\theta]
   =2u^3(S-Y_0)(1-\theta)D_v .
 \tag{EN7}
\]
For precision about exceptional source minors, use the four raw kernel
equations and the two full angular derivatives.  Their selected six-column
system, for \((d,m,A,D,\alpha,\beta)\) with \(C=2\), has
determinant
\[
 -8192\theta u^2vD_u^2D_v^2\,\mathcal D,\qquad
 \mathcal D=Y_0D_u(uY_0+v)(1-uvY_0)
 [S+uv(\theta-1)]R_m(-R_D).
 \tag{EN8}
\]
EN7 follows first from its rational solution wherever this selected
minor is nonzero: the Cramer determinant for the \(m\)-column,
multiplied by \([S+uv(\theta-1)]R_m\), equals this determinant
times \(2u^3(S-Y_0)(1-\theta)D_v\).  This is an exact
polynomial determinant identity.  The equivalent native formula is
\(m=2ae^2c_2/(Mc_3)\), with the definitions in EN13 below:
\(c_2/c_3=(1-\theta)D_v/E\),
\(M=R_m/(Sv)\), \(E=[S+uv(\theta-1)]/S\),
\(a=u(S-Y_0)/v\), \(e=u/S\).  These substitutions give EN7
exactly.  The native contact equations are polynomial identities, so
this argument does not assume the signs being proved.

At the exact positive seed
\[
 (u,v;S,Y_0,\theta)=(1,1;1,\tfrac14,\tfrac12),\qquad
 (d,m,C,A,D,\alpha,\beta)
 =(\tfrac{27}2,4,2,9,\tfrac54,\tfrac{13}4,\tfrac94),
 \tag{EN9}
\]
all six stationary equations hold and the selected minor is nonzero.
Consequently that analytic minor is not identically zero on the
connected contact graph.  Its nonzero set is dense.  Continuity extends
EN7 to every contact on the graph.  At \(\theta=1\), EN7 and \(S,m>0\)
would imply \(R_m=0\), permitting deletion of \(m\) and creating an
exclusive omitted pair, again impossible.  The three continuous
nonzero functions \(Y_0,S-Y_0,\theta-1\) therefore have constant signs,
fixed by EN9:
\[
 S>0,\qquad 0<Y_0<S,\qquad 0<\theta<1.
 \tag{EN10}
\]

The full-parent C1 reflection swaps
\((\lambda,n),(m,C),(A,\alpha),(D,\beta)\), sends \(u\) to \(1/u\),
and sends the state to
\[
 (vS-u(\theta-1),\,-S-uv(\theta-1),\,
                   1-uvY_0,\,uY_0+v).
\]
Its state matrix \(G\) has \(G^TG=D_vI\); the full core and every event
obey \(u^2F_{\sigma_1c}(1/u,v;G\psi)=D_vF_c(u,v;\psi)\).
Its last coordinate is positive.
Applying EN10 to the transformed normalized state proves successively
\(1-uvY_0>0,\ S+uv(\theta-1)>0,\ R_D<0,\ R_m>0\).
Together with EN10 this is precisely EN2.

For a literal third omission \(D=0\) or \(\beta=0\), start at the
given real contact and open that coefficient to \(\varepsilon>0\).
EN4 is still invertible by SE.  The local analytic branch stays at
positive \(d\), has nonnegative event weights, and is physical by SE7.
The already proved inequalities hold on it.  At the limit, EN6 and
SE9 stay strict; the still positive \(A,C,m\) and the same deletion
arguments make the other weak inequalities strict.  EN7 passes to
the limit, and so does the reflection.  Thus EN2 holds on both
boundary faces.  This is an implicit-function limit at a fixed
finite real contact, not an assertion that an optimizer cannot escape
to infinity.  First-order transfer now gives the entire implication
over every real closed field.

\subsection{C2: all signs and a positive omitted-row cofactor}

Use the same positive six-event graph.  Five retained-amplitude walls
are excluded directly by SE10:
\[
 Y=0,\quad uW-Y=0,\quad R_m=0,\quad R_D=0,\quad vY-1=0.
 \tag{EN11}
\]
Deleting \(A,C,m,D,\alpha\), respectively, introduces an exclusive
omitted pair in addition to \(\lambda=\beta=0\).
The wall \(X=0\) deletes \(n\); the resulting source is on the C1
three-consecutive-omission face \(\lambda=n=\beta=0\).
It is physical by SE7, so EN10 gives \(X<0\), a contradiction.
Thus all six retained event rows have fixed signs on the connected
C2 graph.  The exact seed
\[
 (n,m,C,A,D,\alpha)=(\tfrac18,\tfrac38,\tfrac16,
                              \tfrac58,\tfrac38,\tfrac16),\quad d=1,
 \qquad (u,v;W,X,Y,Z)=(1,1;1,-\tfrac14,\tfrac14,1)
\]
has every required positive weight and every stationary equation.
It fixes their signs to the first six signs in EN2.

Here is a direct proof of the omitted-row sign, including \(n=0\).
Let \(P_0\) be this raw C2 matrix with \(n\) set to zero.
Changing \(n\) adds only \(4nD_uD_v\) to its \(XX\) entry.
The cofactor \(B_{24}\) does not contain that column, whereas
\(B_{44}\) changes by \(4nD_uD_v\det(P_0|_{W,Y})\).
Consequently
\[
\begin{split}
 uB_{44}+B_{24}
  ={}&16uv^2D_u^2D_v\,\mathcal F\\
    &+4unD_uD_v\det(P_0|_{W,Y})>0 .
\end{split}\tag{EN12}
\]
Here \(\mathcal F\) is EN6 with \(\beta=0\).  The determinant in
the second line is positive: the \(Z=0\) state block already has
the strictly positive signed-core form SE2.  Thus \(X+u>0\).
The C2 full-parent reflection
\[
 (u,v;W,X,Y,Z)\mapsto(1/v,1/u;Z,X,Y,W),
 \quad (m,D),(C,\alpha),(\lambda,\beta)
\]
has the exact homogeneous identity
\(u^2v^2F_{\sigma_2c}(1/v,1/u;Z,X,Y,W)=F_c(u,v;W,X,Y,Z)\).
The transformed contact
has positive last coordinate \(W\).  Applying the just-proved
\(X'+u'Z'>0\) gives \(W+vX>0\).
This proves EN2 for positive \(n\).  For literal \(n=0\), EN1 is
already a C1 triple face and its signs apply directly; EN12 remains
strict as well.

\subsection{Exact native-coordinate transport and uniqueness}

For C1, reflect if needed so \(0<k=u_c\le1\), and define
\[
\begin{gathered}
 S=W_c,\quad r=Y_c/(kS),\quad
 \theta=1+X_c/k,\quad h=1-\theta,\\
 \eta=v_c/(kS),\quad \zeta=Sv_c/k,\quad
 f=\sqrt\eta,\quad g=\sqrt\zeta,\quad e=k/S=kf/g,\\
 a=(1-r)/\eta,\quad t=\theta/\zeta,\quad x=k^2,\\
 R=1+xr,\quad T=1+xh,\quad M=R-t,\quad J=T-a,\\
 E=1-x\eta h,\qquad Q=1-x\zeta r .
\end{gathered}\tag{EN13}
\]
By EN2, \(0<r,\theta<1\), all \(\eta,\zeta,a,t\) are positive,
and
\[
 R_m=k\zeta M>0,\quad R_D=-v_cJ<0,\quad
 W_c+v_cX_c=SE>0,\quad 1-v_cY_c=Q>0.
\]
Thus \(M,J,E,Q>0\), and \(0<a,t<2\).
The selected determinant EN8 is now strictly nonzero.  The native
seven-coquotient formulas of \emph{native\_formulas.tex}, together
with the source coefficients in SD1 of \emph{source\_definitions.tex},
satisfy exactly the six stationary equations at this same contact.
Uniqueness therefore identifies them with the original normalized
source, including \(D=0\) or \(\beta=0\).  No asymptotic hypothesis
is used for these polynomial identities.

For C2, normalize \(d=1\), and set
\[
\begin{gathered}
 c=u_c,\quad a=v_c/u_c,\quad b=aW_c,\quad r=-X_c/u_c,\quad
 t=(u_cW_c-Y_c)/v_c,\quad z=c^2,\\
 h=b-a^2t=aY_c/u_c,\quad
 g=(1+z)b-(1-r+a^2zt)=R_m/u_c,\quad
 j=1+zr-t=-R_D/v_c .
\end{gathered}\tag{EN14}
\]
The complete strict native domain follows:
\[
 a,b,c,t>0,\quad 0<r<1,\quad h,g,j>0,\quad
 1-zh=1-v_cY_c>0,\quad b-a^2zr=a(W_c+v_cX_c)>0.
 \tag{EN15}
\]
The actual six-column stationary determinant is, in the unweighted
event convention of the retained determinant proof,
\[
 -64ac^6rgt\,hj(1-zh)(1+z)^2(1+a^2z)^2\Delta,
\]
\[
 \Delta=1+(1+z)b+z(1+z)b^2+2z(1+z)br+z^2r^2+2zr
        +a^2z\{(1+z)b+1+zr(2-r)\}\ge1 .
 \tag{EN16}
\]
That body provides its explicit reduction to a three-by-three
determinant.  The six rational coefficients in
\emph{source\_formulas.tex} solve these six equations at the
contact EN14.  EN16 makes the solution unique, so they are exactly
the original source coefficients.  This remains valid at \(n=0\);
one does not delete its column from the square system.

All gauge actions, positivity statements and source equations used
above are finite polynomial or rational ordered-field statements
with denominators explicitly nonzero.  The real connected-graph
proof therefore transfers as one fixed first-order implication.
Together, EN1--EN16 discharge the native sign and positive-gap
interface for both families in SE15.  They do not prove an
all-tests tensor assertion; that assertion is supplied by the
subsequent exhaustive source partition and its actual terminal bodies.

